\documentclass[11pt]{book}
\usepackage[utf8]{inputenc}
\usepackage{geometry}
\geometry{a5paper, margin=1.5cm}
\usepackage{ebgaramond}
\usepackage{epigraph}
\usepackage{verse}
\usepackage{microtype}
\setlength{\epigraphwidth}{0.7\textwidth}
\setlength{\epigraphrule}{0pt}

\begin{document}

\chapter*{The Justicar's Judgment}
\addcontentsline{toc}{chapter}{The Justicar's Judgment}

\vspace{2cm}
\begin{center}
\textit{A tale of Viterra, told by a clerk who sharpened the quills}
\end{center}
\vspace{1cm}

\epigraph{The law does not forget. Neither does the son who was not ransomed.}{--- Queen's Justicar Elodie Venn, from the bench}

%--- Prologue ---------------------------------------------------------
\section*{Prologue – The Hedge‑Court}
The evening mist clung to Dusana's hedge‑court like a shy lover, curling around the low stone benches and the half‑filled mug of small beer that steamed faintly in the chill. Dusana set the cup down, the amber liquid catching the last glint of daylight, and laid before her a vellum sheet black‑waxed with the crown of Viterra and the broken column of Vhasia. The smell of iron and forgotten oaths lingered in the air.

\bigskip
\textit{You are commanded to hear the case of Baron Aldric of the Scar…}

\bigskip
Her hand did not tremble.

She remembered the boy who had once been Sir Aldric of the White Moor—a laughing knight whose sword sang as he rode, whose grin could coax a smile from a storm. She had borne his shield at the Ford of Ashes, watched him trade a burnt tower for a handful of serfs with the same eager barter he used over a horse, and had seen him abandon Queen Veralyn’s mother at the Battle of the Three Bridges, fling his banner to the Pretender’s side, then swear fealty again when the Pretender fell. At the first inquiry she had spoken for him, believing in second chances as if they were coins she could spend without limit.

Now the queen’s justice asked her to weigh the man who had once been her mentor against the scales of law and memory.

%--- Scene 3 – Chapel‑Burner (Dusk) -------------------------------------------
\section*{Scene 3 – Cross‑Examination: The Chapel‑Burner}
\smallskip
\textbf{Call the witness: Brother Tomas, keeper of Saint Mirelle’s Chapel.}

The light was already slipping toward dusk when Tomas shuffled forward, his robe frayed at the hem, eyes haunted by the memory of flame. He took the oath, his voice a thin reed trembling in the gathering gloom.

\begin{description}
\item[Elodie (softly):] ``Brother Tomas, tell the court what you saw on the night the chapel burned.''
\item[Tomas (voice trembling):] ``My lord, I was at vespers when the torches came—smoke already curled from the nave like a warning. Men in the Baron’s colours shattered the door, shouting that rebels hid within. They thrust flaming brands into the straw, and the fire rose like a hungry beast, licking the beams and devouring the rafters. I fought my way through the heat, pulling two children from the pews before the ceiling gave way; twelve souls were lost that night, sir.’’
\item[Alcric’s counsel (leaning forward, voice oily):] ``And you, brother, were you not once a tenant of the Baron’s lands? Might your testimony be coloured by a grudge?’’
\item[Tomas (flushing, hands clasping the rosary tighter):] ``I bear no grudge, only the memory of the innocent who perished. I speak for them, not for myself.’’
\end{description}

The court murmured, the weight of those twelve lives settling like ash on the shoulders of those who heard, the fading light casting long, somber shadows across the stone floor.

\section*{Scene 4 The Wetnurse’s Hook}
\smallskip
\textbf{Call the witness: Bryn, former wetnurse to Lady Margot, first wife of the accused.}

The gallery parted slowly. A woman in grey, leaning on a knotted staff, made her way to the bar. Her left hand ended in a brass hook—rough, forged from a farmer’s rasp, still scarred by the fire where it had been shaped. Her face was a map of old grief, the skin drawn tight over sharp bones. She seemed older than the stone of the courthouse, yet her eyes, when they found Aldric, burned with a light that had not dimmed in forty years.

A murmur ran through the onlookers. Those who knew the grey robes, the brass hook, the rumors of the Cord, shifted on their benches. Some touched their own wrists where a red thread might once have been. Others stared at the floor.

\begin{description}
\item[Elodie:] ``State your name and relation to the accused.''
\item[Bryn (voice thin as winter ice, but clear):] ``I am Bryn. I served as wetnurse to his first wife, Lady Margot, and to their daughter, Eleanor. I was midwife to her when she died.''
\item[Aldric’s counsel (rising, a sneer in his tone):] ``A servant’s testimony from forty years past? The woman is ancient, her memory clouded by age and spite. Moreover, the Baronet held pit and gallows on his own land. Any discipline he administered there is not subject to this court’s review.''
\item[Elodie (after a long pause, her gavel resting on the bench):] ``The court will hear the witness, though the baron’s right to enforce law on his holding is noted. Proceed.''
\end{description}

Bryn’s good hand reached into her grey robe and drew out a silver locket, tarnished but still whole. She held it up.

\begin{quote}
``Margot gave me this on her deathbed. ‘For the child,’ she said. ‘The one who still breathes.’ Aldric never asked for it. He never visited her grave. I left his service in disgust. Three days later, his men found me at a waystation. He accused me of thievery – of stealing this locket, which he had never once mentioned while his wife lived. He did not ask. He just swung. He laughed and said the hand was enough; I could keep the ‘trinket.’ ''
\end{quote}

She raised her brass hook toward the bench. The light caught its rough surface.

\begin{quote}
``This was forged from a farmer’s rasp, at a forge that no longer stands. The hand that fed his children, that held his dying wife’s head – he took it for a trinket he had forgotten existed. I left the Cord behind. But I do not forget. Neither does the law, if it has the courage to remember.''
\end{quote}

A heavier murmur passed through the gallery. Some nodded; others looked away. Aldric’s counsel objected again, citing the baron’s right of pit and gallows. Elodie, after a whispered consultation with the Codex clerk, ruled the testimony admissible only as character evidence; the baron’s disciplinary authority on his own land could not be challenged by the Queen’s court. The substance of the accusation – the loss of Bryn’s hand – was dismissed as beyond the court’s jurisdiction.

But the locket remained on the evidence table. Bryn did not weep. She lowered her hook and limped back into the gallery, her eyes never leaving Aldric’s face. He did not look at her. He could not. But he smirked.

%--- Scene 4 – Iron Tithe (Dawn) ---------------------------------------------
\section*{Scene 5 – Cross‑Examination: The Iron Tithe}
\smallskip
\textbf{Call the witness: Mirella, widow of the miller of the burnt tower.}

Dawn had just broken, a thin veil of mist clinging to the fields as Mirella stepped forward, her hands calloused from years of grinding grain, her eyes hard as flint against the pale light.

\begin{description}
\item[Elodie:] ``Mirella, describe the levy the Baron imposed on your household.''
\item[Mirella (steady, voice low but resonant):] ``He arrived with his men at first light, declaring that protection required half our harvest. We gave it, trusting his word, believing the promise of safety. Then he carted the grain to market, selling it at full price while we watched our stores dwindle. My children grew thin; my husband fell ill from the strain, his cough echoing in the empty barn.’’
\item[Alcric’s counsel (smirking, voice smooth as polished stone):] ``Was it not the Baron’s right to demand tribute in times of war? Did you not benefit from the peace his forces brought?’’
\item[Mirella (voice rising, a flash of anger breaking through the dawn calm):] ``Peace bought with our blood is no peace at all. He took our sustenance to fund a private militia that terrorised the very folk he swore to protect. Each sack of grain felt like a blade against our throats.’’
\end{description}

A ripple of disapproval passed through the nobles; the crisp morning air seemed to carry the metallic tang of greed, sharpening the sense that justice, though slow, was finally stirring.

\section*{Scene 6 – The Son Who Was Not Ransomed}
\smallskip
\textbf{In the year of the Broken Crown, after the Battle of the Red Ditch, your son—Ser Roran, then seventeen—was captured by the forces of the Duke of Adria. The customary ransom for a knight’s eldest son was three hundred silver marks. The Duke’s herald offered you a discount of fifty marks if you paid within the week.}

Aldric’s smile thinned to a razor’s edge.

\begin{quote}
``I recall.''
\end{quote}

\begin{quote}
``You said, and I quote the herald’s sworn testimony: ‘I have more sons and can make further. The boy was weak. Let the Duke keep him.’''
\end{quote}

A rustle stirred the gallery. In the shadows sat a lean, grey‑faced man—Ser Roran, now thirty‑seven, his left hand sheathed in a silvered gauntlet, his right resting on the pommel of a sword that had seen twenty campaigns.

Elodie continued, her voice low but resonant, each syllable a thread pulling the tapestry tighter.

\begin{quote}
``The Duke did not keep him. He threw the boy into a debtor’s cell, expecting you to relent. You did not. For three years Roran rotted. Then a stranger—once a hedge‑knight under your banner—paid the ransom with every coin he had, died of fever the following winter, and gave Roran his name and his sword.''
\end{quote}

She paused, letting the image of the boy’s lonely cell settle like frost on stone.

\begin{quote}
``He did not return to you. He fought in the eastern wars. He won the Queen’s own pennon at the Siege of the Horn. He knighted himself with a broken blade and the blessing of a battlefield priest. He has not spoken your name in seventeen years.''
\end{quote}

Aldric’s jaw tightened, the muscle jumping like a trapped hare.

\begin{quote}
``The boy survived. Stronger for it.''
\end{quote}

Ser Roran rose. The hall fell silent as he walked to the bar, placed his gauntleted hand on the rail, and spoke, his voice barely above a whisper yet carrying the weight of years.

\begin{quote}
``Father, do you know what I remember most? Not the cell. Not the chains. The moment I heard you say it. I was hiding behind the herald’s tent. I heard you laugh.''
\end{quote}

\section*{Scene 7 – The Justicar’s Judgment}
The law of Viterra stood ancient and immutable: treason meant death; oath‑breaking stripped lands and titles; refusing a son’s ransom was not a crime—only cruelty.

Elodie spent three hours upon the bench, the massive \textit{Codex of Precedents} hauled by two clerks like a tombstone. She asked for water, for silence, for the space to let the truth settle like dust after a storm.

When she rose, her voice was level, each word a stone placed carefully upon the scales of justice.

\begin{quote}
``Baron Aldric of the Scar, on the first count—treason by repeated oath‑breaking—I find you guilty. The sentence is forfeiture of all lands, titles, and pensions. You will be stripped of your banner and your ring. You will live as a hedge‑knight if any will take you.''
\end{quote}

Aldric’s face remained unmoved; he had expected this.

\begin{quote}
``On the second count—refusal of ransom and abandonment of kin—the law is silent. Viterra does not compel a father to love. But the law does possess a clause for debts of the heart: the \textit{Right of Acknowledgment}. No one has invoked it in two hundred years.''
\end{quote}

She turned her gaze to Ser Roran.

\begin{quote}
``Ser Roran, son of Aldric, you hold the right to demand that your father speak one truth aloud, in this court, under witness. Should he refuse, he is declared \textit{anima non grata}—his soul outlawed from all Vritann burial grounds. His body will be interred at a crossroads with a stake through his chest. The law permits it; the Church condemns it. Yet the law is the law.''
\end{quote}

Aldric laughed—a short, bitter bark.

\begin{quote}
``A truth? I have spoken a thousand truths. I have said a man must cut what he cannot carry. I have said sons are arrows—you loose them and you do not weep. That is my truth.''
\end{quote}

Ser Roran stepped forward, his voice soft as falling silk yet unyielding.

\begin{quote}
``No. Your truth is that you were afraid. Not of the Duke. Of me. I was stronger than you at seventeen. You saw it. You could not bear to pay my ransom because then you would owe me, and you have never owed anyone in your life.''
\end{quote}

Aldric opened his mouth, closed it. The hall held its breath, the silence thick as the velvet drapes.

\begin{quote}
``I am not asking you to love me,'' Roran said. ``I am asking you to say it. Just once. Say: ‘I was afraid.’''
\end{quote}

Aldric stared at his son. The mask of the turncloak—the smile, the shrug, the armor of indifference—cracked like thin ice under a spring thaw.

\begin{quote}
``You killed my wife,'' he growled. "You wear her face."
\end{quote}

Roran stared at him, hard, waiting. \textit{"Liar."} The accusation was written in his eyes.

\begin{quote}
    ``I was afraid,'' he finally whispered.
\end{quote}

The sound was so small the clerks leaned forward, straining to catch it.

\begin{quote}
``Of what?'' Elodie prompted.
\end{quote}

\begin{quote}
``Of him. Of what he would become. He would have been better than me. And I… I could not stand to watch.''
\end{quote}

Ser Roran turned away, his steps silent upon the stone, the silver gauntlet clicking a faint, steady rhythm. He did not weep; the tears had long since dried into salt upon his cheeks.

Elodie lifted her gavel.

\begin{quote}
``The Right of Acknowledgment is satisfied. The court is adjourned. Baron Aldric—former Baron Aldric—you are free to go. But I advise you to leave Viterra by sunset. The road is long, and the winter is coming.''
\end{quote}

\section*{Epilogue – Aftermath}
Three winters later Aldric was found dead in a ditch outside Thepyrgos, near the Black River, reeking of wine. His lay body unclaimed, circled in salt. The crossroads stake was unnecessary—he had outlived his own burial rights. Scratched into his chest near a circular wound were the words \textit{anima non grata}

Ser Roran entered the Queen’s household as a knight, wearing the silvered gauntlet over his left hand—not to conceal the scar, but to remind himself that a stranger’s coin had bought his life. He never took ransom. Once a captured lord offered him five hundred marks for his freedom.

\begin{quote}
``Pay it to the hedge‑knight’s widow,'' he said. ``I have no need for gold. I had a father once. That was expensive enough.''
\end{quote}

When Queen Veralyn learned of the verdict, she sat in silence for a long breath, then summoned the \textit{Codex of Precedents} and, in her own hand, added a marginal note that would become law:

\begin{quote}
``The Right of Acknowledgment may be invoked by any child who was not ransomed, not visited, not named in a will. The truth, once spoken, is payment in full. No further interest shall accrue.''
\end{quote}

The scribes copied the clause into every new edition. It remains law to this day.

\bigskip
\noindent\textit{The ink is dry. The gavel is still.}

\end{document}